%% Packages
\documentclass[12pt,a4paper]{report}
\usepackage[T1]{fontenc}
\usepackage[utf8]{inputenc}
\usepackage{minitoc}
\usepackage[pdftex]{graphicx}
\usepackage[a4paper,left=2cm,right=2cm,top=2cm,bottom=2cm]{geometry}
\usepackage[frenchb]{babel}
\usepackage{libertine}
\setcounter{tocdepth}{1}

\setcounter{minitocdepth}{3}

%% Document

\begin{document}

\begin{titlepage}
	\begin{center}

		\includegraphics[scale=1.5]{garde.jpg}~\\[1.5cm]
		\textsc{\LARGE RAPPORT DE PROJET BUREAU D'ETUDES}\\[2cm]	
	
	\end{center}
\end{titlepage}



   \dominitoc
   \tableofcontents
   \chapter{Le projet}
      \section{Sujet}
      	%Le  projet  consiste  à  déterminer  le  temps  de  demi-vie  du  polonium 212 qui  se désintègre en émettant une particule α dont l’énergie est comprise entre 8 MeV et 9 MeV. Dans un premier temps, vous utiliserez le programme qui vous donnera une valeur de la probabilité de sortie de la particule α en dehors du noyau et à l’aide d’un modèle que vous aurez élaboré, vous en déduirez un ordre de grandeur du temps de demi-vie. Dans un second temps vous comparerez la valeur ainsi obtenue aux valeurs expérimentales et celle déterminée à l’aide de la formule de Gamow.
      \section{Lien avec le cours}
   \chapter{Réalisation}
   		\section{Première étape}
   		\section{Deuxième étape}
\end{document}